\documentclass[10pt,twocolumn,letterpaper]{article}

\usepackage{iccv}
\usepackage{times}
\usepackage{epsfig}
\usepackage{graphicx}
\usepackage{amsmath}
\usepackage{amssymb}

\def\eg{{\emph{e.g.}}}
\def\ie{{\emph{i.e.}}}
\def\etc{{\emph{etc.}}}
\def\etal{{\emph{et al.}}}
% Include other packages here, before hyperref.
\usepackage{graphicx}
\usepackage{amsmath}
\usepackage{amssymb}
\usepackage{booktabs}
\usepackage{pifont} % colors
\usepackage{amsmath}
\usepackage{amssymb}
\usepackage{graphicx}
\usepackage{array}
\usepackage{comment}
\usepackage{epsfig}
\usepackage{times}
\usepackage{comment}
\usepackage{url}
\usepackage{subfig}
\usepackage{times}
\usepackage{algorithm,algorithmic,xspace}
\usepackage{verbatim}
\usepackage{multirow}
\usepackage{colortbl}
\usepackage{bbding}
\usepackage{pifont}
\usepackage{wasysym}
\usepackage{amssymb}
\usepackage[table]{xcolor}
\usepackage{breqn}
\usepackage{graphicx,animate}
\usepackage[normalem]{ulem}
\usepackage{tikz}
\newcommand*{\circled}[1]{\lower.7ex\hbox{\tikz\draw (0pt, 0pt)%
    circle (.4em) node {\makebox[1em][c]{\small #1}};}}

\newcommand{\Tref}[1]{Table~\ref{#1}}
\newcommand{\Eref}[1]{Equation~(\ref{#1})}
\newcommand{\Fref}[1]{Figure~\ref{#1}}
\newcommand{\Sref}[1]{Section~\ref{#1}}
\newcommand{\Aref}[1]{Algorithm~\ref{#1}}

\def\etal{\emph{ et al.}}
\def\ie{\emph{i.e.}}
\def\eg{\emph{e.g.}}
\def\httilde{\mbox{\tt\raisebox{-.5ex}{\symbol{126}}}}
\def\c{\mathbf{c}}
\def\F{{\mathcal{F}}}
\def\eg{{\emph{e.g.}}}
\def\ie{{\emph{i.e.}}}
\def\etc{{\emph{etc.}}}
\def\etal{{\emph{et al.}}}
\def\L{{\mathbf{L}}}
\def\I{{\mathbf{I}}}
\def\B{{\mathbf{B}}}
\def\M{{\mathbf{M}}}
\def\Ch{\hat{\mathbf{C}}}
\def\C{\mathbf{C}}

\def\z{{\mathbf{z}}}
\def\x{{\mathbf{x}}}
\def\G{{\mathbf{\mathcal{G}}}}
\def\ray{{\mathcal{R}}}
\def\d{{\mathbf{d}}}
\def\P{{\mathbf{P}}}
\def\o{{\mathbf{o}}}
\def\loss{{\mathcal{L}}}
\def\d{{\mathbf{d}}}
\def\K{\mathcal{K}}
%\def\Pi{{\mathbf{P}_{i}}}
\def\Pi{{\mathbf{M}_{i}}}
\def\r{{\mathbf{r}}}
\def\T{{\mathbf{T}}}
\def\H{{\mathbf{H}}}
\def\W{{\mathbf{W}}}
\def\G{{\mathbf{G}}}
\def\bb{\mathbf{b}}
\def\cta{\mathbf{\theta}}
\def\nerf{NeRF~}
\def\CopyNeRF{CopyRNeRF~}

\definecolor{Auqamarin}{gray}{0.9}
\definecolor{LightCyan}{rgb}{0.88,1,1}
\definecolor{myGreen}{rgb}{0, .9, .6}
\newcommand{\renjie}[1]{\textcolor{red}{{#1}}}
\newcommand{\boxin}[1]{\textcolor{magenta}{{[Boxin: #1]}}}
\newcommand{\ziyuan}[1]{\textcolor{blue}{{#1}}}
\definecolor{americanrose}{rgb}{1.0, 0.01, 0.24}
\newcommand{\qing}[1]{\textbf{\textcolor{purple}{qing: #1}}}

\newcommand{\argmax}{\operatornamewithlimits{argmax}}
\newcommand{\argmin}{\operatornamewithlimits{argmin}}
\newcommand{\revision}[1]{\textcolor{red}{{#1}}}
\makeatletter
\newcommand*{\rom}[1]{\expandafter\@slowromancap\romannumeral #1@}
\DeclareRobustCommand\onedot{\futurelet\@let@token\@onedot}
\def\@onedot{\ifx\@let@token.\else.\null\fi\xspace}

% \newfloat{figtab}{htb}{fgtb}
\makeatletter
  \newcommand\figcaption{\def\@captype{figure}\caption}
  \newcommand\tabcaption{\def\@captype{table}\caption}
\makeatother
% Include other packages here, before hyperref.

% If you comment hyperref and then uncomment it, you should delete
% egpaper.aux before re-running latex.  (Or just hit 'q' on the first latex
% run, let it finish, and you should be clear).
\usepackage[breaklinks=true,bookmarks=false]{hyperref}

\usepackage[capitalize]{cleveref}
\crefname{section}{Sec.}{Secs.}
\Crefname{section}{Section}{Sections}
\Crefname{table}{Table}{Tables}
\crefname{table}{Tab.}{Tabs.}
\iccvfinalcopy % *** Uncomment this line for the final submission

\def\iccvPaperID{6271} % *** Enter the ICCV Paper ID here
\def\httilde{\mbox{\tt\raisebox{-.5ex}{\symbol{126}}}}

% Pages are numbered in submission mode, and unnumbered in camera-ready
\ificcvfinal\pagestyle{empty}\fi

\begin{document}

%%%%%%%%% TITLE
\title{CopyRNeRF: Protecting the CopyRight of Neural Radiance Fields}


\author{Ziyuan Luo$^{1,2}$\quad Qing Guo$^{3}$ \quad Ka Chun Cheung$^{2, 4}$ \quad Simon See$^{2}$ \quad Renjie Wan$^{1}$\thanks{Corresponding author. This work was done at Renjie's Research Group at the Department of Computer Science of Hong Kong Baptist University.}\\
$^{1}$\small{Department of Computer Science, Hong Kong Baptist University}\\
$^{2}$\small{NVIDIA AI Technology Center, NVIDIA}\\
$^{3}$\small{IHPC and CFAR, Agency for Science, Technology and Research, Singapore}\\
$^{4}$\small{Department of Mathematics, Hong Kong Baptist University}\\
{\tt\small ziyuanluo@life.hkbu.edu.hk, guo_qing@cfar.a-star.edu.sg, \{chcheung, ssee\}@nvidia.com,}\\ {\tt\small renjiewan@hkbu.edu.hk}
}

%\author{Ziyuan Luo\\
%Institution1\\
%Institution1 address\\
%{\tt\small firstauthor@i1.org}
% For a paper whose authors are all at the same institution,
% omit the following lines up until the closing ``}''.
% Additional authors and addresses can be added with ``\and'',
% just like the second author.
% To save space, use either the email address or home page, not both
%\and
%Qing Guo\\
%Institution2\\
%First line of institution2 address\\
%{\tt\small secondauthor@i2.org}
%\and
%Ka Chun Cheung\\
%Institution2\\
%First line of institution2 address\\
%{\tt\small secondauthor@i2.org}
%\and
%Simon See\\
%Institution2\\
%First line of institution2 address\\
%{\tt\small secondauthor@i2.org}
%\and
%Renjie Wan\\
%Institution2\\
%First line of institution2 address\\
%{\tt\small secondauthor@i2.org}
%}

\maketitle
% Remove page # from the first page of camera-ready.
\ificcvfinal\thispagestyle{empty}\fi

%%%%%%%%% ABSTRACT
\begin{abstract}
Neural Radiance Fields (NeRF) have the potential to be a major representation of media. Since training a NeRF has never been an easy task, the protection of its model copyright should be a priority. In this paper, by analyzing the pros and cons of possible copyright protection solutions, we propose to protect the copyright of NeRF models by replacing the original color representation in NeRF with a watermarked color representation. Then, a distortion-resistant rendering scheme is designed to guarantee robust message extraction in 2D renderings of NeRF. Our proposed method can directly protect the copyright of NeRF models while maintaining high rendering quality and bit accuracy when compared among optional solutions. Project page: \href{https://luo-ziyuan.github.io/copyrnerf}{https://luo-ziyuan.github.io/copyrnerf}.

%introduce an approach to directly watermark the 

%we propose protecting the ownership of such implicit 3D representations by ensuring that any 2D images generated from them are protected with an invisible copyright message. By manipulating the color factor used for rendering, our method can effectively embed copyright messages into NeRF's implicit representation and ensure high-quality reconstruction. A distortion-resistant rendering is designed to ensure the robust message extraction in rendered 2D images even with a conventional CNN-based message extractor. Compared with existing methods, our method can guarantee invisible message embedding and robust message extraction on the evaluation dataset.%
%\qing{Is that possible to highlight the main challenges and differences to existing 2D watermarking methods?}
%
\end{abstract}

%%%%%%%%% BODY TEXT
\section{Introduction}
\label{sec:intro}




%he bit accuracy of watermark extraction.
\begin{figure}[t]
  \centering
%   \animategraphics[loop,autoplay,width=\linewidth]{2}{figs/figure1-}{0}{1}
  \includegraphics[width=\linewidth]{figs/fig_1.pdf}
  \caption{When NeRF models are stolen (\protect\circled{1}) by maclicious users, \CopyNeRF can help to claim model ownership by transmitting copyright messages embedded in models to rendering samples (\protect\circled{2}). We show some comparisons with HiDDeN~\cite{zhu2018hidden} + NeRF~\cite{nerf2020}, and NeRF~\cite{nerf2020} with messages. PSNR/Bit Accuracy is shown below each example.}
  \label{fig:short}
      \vspace{-15pt}
\end{figure}

%an implicit 3D representation can be simply constructed via a set of 2D images using NeRF~\cite{nerf2020}, and then directly employed to synthesize images with novel viewpoints. Despite its convenience Rarely do people question the importance of intellectual property. 

Though Neural Radiance Fields (NeRF)~\cite{nerf2020} have the potential to be the mainstream for the representation of digital media, training a NeRF model has never been an easy task. If a NeRF model is stolen by malicious users, \textit{how can we identify its intellectual property?}

As with any digital asset (\eg, 3D model, video, or image), copyright can be secured by embedding copyright messages into assest, aka digital watermarking, and NeRF models are no exception. 
An intuitive solution is to directly watermark rendered samples using an off-the-shelf watermarking approach (\eg, HiDDeN~\cite{zhu2018hidden} and MBRS~\cite{jia2021mbrs}). However, this only protects the copyright of rendered samples, leaving the core model unprotected. If the core model has been stolen, malicious users may render new samples using different rendering strategies, leaving no room for external watermarking expected by model creators. Besides, without considering factors necessary for rendering during watermarking, directly watermarking rendered samples may leave easily detectable trace on areas with low geometry values.

%The digital watermarking of 3D models can be conducted by embeding messages directly into 3D structures. 
The copyright messages are usually embedded into 3D structure (\eg, meshes) for explicit 3D models~\cite{Yoo_2022_CVPR}. Since such structures are all implicitly encoded into the weights of multilayer perceptrons (MLP) for NeRF, its copyright protection should be conducted by watermarking model weights. As the infomration encoded by NeRF can only be accessed via 2D renderings of protected models, two common standards should be considered during the watermark extraction on rendered samples~\cite{ahmadi2020redmark, jing2021hinet, Yang_2021_ICCV, zhang2020udh}: 1) \textbf{invisibility}, which requires that no serious visual distoration are caused by embedded messages, and 2) \textbf{robustness}, which ensures robust message extraction even when various distortions are encountered. 


%even if the rendered image is severely distorted or even different rendering approaches are employed. 

%An intuitive solution is to directly watermark rendered samples using an off-the-shelf watermarking approach (\eg, HiDDeN~\cite{zhu2018hidden}). However, this only protects the copyright of rendered samples, leaving the core model unprotected. Besides, without considering the geometry, directly watermarking rendered samples may leave easily detectable trace. Thus, it is pivotal to focus on how to watermark NeRF model while maintaining two common standards through rendered samples.

%One pioneer method has been proposed for the watermarking of explicit 3D models~\cite{yoo2022deep}. By directly embedding messages into 3D meshes, their approach guarantess the standard for invisibility and robustness in rendered samples. However, it is impossible to find such explicit 3D meshes to store secret messages in NeRF, since all information are implicitly encoded into the weights of a connectionist representation (\eg, MLP)~\cite{nerf2020,zhi2021place}. 


One option is to create a NeRF model using watermarked images, while the popular invisible watermarks on 2D images cannot be effectively transmitted into NeRF models. As outlined in~\Fref{fig:short} (HiDDeN~\cite{zhu2018hidden} + NeRF~\cite{nerf2020}), though the rendered results are of high quality, the secret messages cannot be robustly extracted. We can also directly concatenate secret messages with input coordinates, which produces higher bit accuracy (NeRF with message in~\Fref{fig:short}). However, the lower PSNR values of rendered samples indicate that there is an obvious visual distortion, which violates the standard for invisibility. 
%Thus, the problem becomes how to embed messages into NeRF models, while maintaining the two standards in rendered samples.

%However, invisible watermarks are typically embedded into 2D images without obviously affecting color and geometry. Since NeRF only relies on a simple MLP to estimate color and geometry for volume rendering, it is incapable of capturing such invisible changes during scene modeling. As outlined in~\Fref{fig:short} (HiDDeN~\cite{zhu2018hidden} + NeRF~\cite{nerf2020}), though the rendered results are still of high quality, the secret messages cannot be robustly extracted. Besides, we can also directly watermark 3D models. For instance, several recently proposed methods~\cite{yoo2022deep} have chosen to directly embed information in 3D meshes and access them from 2D renderings. However, it is impossible to find such explicit 3D meshes to store secret messages in NeRF, since all information are implicitly encoded into the weights of a connectionist representation (\eg, MLP)~\cite{nerf2020,zhi2021place}.


%Another approach is to watermark 3D models directly. Traditionally, this kind of technologies can be divided into two categories. One is that messages are embedded and extracted both in 3D space. However, this method will fail in scenes where 3D models cannot be accessed directly, such as movies and games. In these cases, messages can only be retrieved by 2D images rendered by 3D models. The other is embedding information in texture and vertex of 3D mesh and extracting information from rendered 2D image. However, this method cannot generalize to complex 3D models limited to its renderer, and NeRF model cannot be used as input. To the best of our knowledge, there is no relevant research on applying watermark technology to NeRF at present.

%Considering that the model only needs to be supervised in 2D space during the training process of NeRF, and the renderer is differentiable, we propose a watermarking framework in NeRF. The proposed method can embed the watermark into the implicit representation of NeRF, and extract information from 2D images rendered from any views without changing the perceptual content.For example, as shown in~\Fref{fig:short}, the visual quality still looks promising when we follow the strategy proposed in~\cite{yu2021pixelnerf} to concatenate the message with the input coordinates, while the message extraction accuracy is low.

%Considering the implicit property of NeRF models, we can also embed secret messages into network parameters. Many strategies~\cite{yu2021pixelnerf, chen2021fast} have been proposed to incorporate certain information into NeRF. For example, PixelNeRF~\cite{yu2021pixelnerf} concatenates additional features with input coordinates. However, as shown in~\Fref{fig:short} (NeRF with messages), concatenating messages with input coordinates may affect the reconstruction quality.  Moreover, the pixel-wise rendering in vanilla NeRF is unable to effectively consider distortions for 2D images during training, since most distortions (\eg, blurring and cropping) are only meaningful in a certain area. That may undermine the robustness of message extraction, when the rendered images encounter distortions. In addition, the novel views synthesized from a NeRF model may be with different viewpoints to the training views. How can we maintain robust message extraction even under new synthesized viewing perspectives?

%How can we ensure the robust message extraction when encountering rendered images with totally different viewpoints? 

%Moreover, NeRF and its variants actually build a local correlation by mapping a spatial point in 3D to an image pixel in 2D. 
%\qing{I do not understand the following issues, what are the global distorations? How does it affect?}
%Any information embedded via such local mapping may easily be violated when encountering global distortions (\eg, cropping and blurring). In order to comply with the requirement of invisibility and robustness in digital watermarking, the usual method used by ordinary NeRF and its variants cannot be directly utilized.

%One alternative way to manipulate the network parameters is to directly use secrete messages as a component of the input, and further use a message extractor to access secret messages from the rendered 2D images. However, by using a secret message directly as the input, it significantly affects the geometry and color of rendered images, which makes the observer easily feel the existence of additional information (as shown in Fig. xx). Moreover, NeRF and its variants actually build a local correlation by mapping a spatial point in 3D to an image pixel in 2D. Any information embedded via such local mapping may easily be violated when encountering global distortion (\eg, cropping and blurring). Considering the importance of invisibility and robustness in digital watermarking, our framework should be able to ensure the two aspects in our framework.

%One alternative way is to directly use secrete messages as a component of the input, and further use a message extractor to access secret messages from the rendered 2D images. However, directly embedding a secret message into the coordinates may significantly affect the geometry and color of rendered 2D images, which makes the observer easily feel the change of contents. Moreover, by querying hundreds of sample points along a ray, NeRF and its variants actually builds a mapping from a spatial position in 3D to an image pixel in 2D. Any message directly embedding along with such a local mapping cannot be effectively accessed by a message extractor designed for the whole image. }

%actually construct a mapping from a spatial position in 3D to an image pixel in 2D, while


%However, there are two challenges to leverage NeRF as the backbone of 3D scene watermarking. Firstly, directly embedding message to coordinates to predict the  watermarked opacity and radiance color will affect the geometry of 3D scenes, which means the non-scene information will generate space inconsistent results
%across different views. This is because the images with embedded messages lack 3D information. Secondly, NeRF needs to query hundreds of sample points along the ray to render a single pixel. However, the extraction of information is usually based on the global information of the image, and pixel-based rendering will go against this.

%\renjie{If hiding messages within explicit 3D structure or the implicit mapping are all not available, we propose to embed messages by directly manipulating the intermediate factors during NeRF-based scene representation? Two intermediate factors can be obtained the NeRF-based representation: color and opacity, where the former one is view-dependent emitted radiance, and the latter one denotes differential likelihood of a ray hitting a particle~\cite{https://doi.org/10.48550/arxiv.2209.02417}.

%In this work, we embed the secret messages into the color component due to a simple reason: Geometry information is difficult to be directly accessed from 2D images. \textbf{Some sentences are needed to say why such strategy can preserve information.}}, as 2D projection images lack of 3D spatial information\cite{Ge_2017_CVPR, 8338122}.

% \cite{bas2019does, Chen_2019_ICCV, Qin_Wang_Lu_2019, 8910034, gong2016three, 8338122}the components (geometry and color) used for volume rendering in NeRF. Though hiding in the geometry is also an option, the geometry in 2D images is with ambiguity~\cite{bas2019does}, which may undermine the robustness of message extraction. 
 
Though invisibility is important for a watermarking system, the higher demand for robustness makes watermarking unqiue~\cite{zhu2018hidden}. Thus, in addition to invisibility, we focus on a more robust protection of NeRF models. As opposed to embedding messages into the entire models as in the above settings, we create a \textit{watermarked color representation} for rendering based on a subset of models, as displayed in~\Fref{fig:framework}. By keeping the base representation unchanged, this approach can produce rendering samples with invisible watermarks. By incorporating spatial information into the watermarked color representation, the embedded messages can remain consistent across different viewpoints rendered from NeRF models. We further strengthen the robustness of watermark extraction by using \textit{distortion-resistant rendering} during model optimization. A distortion layer is designed to ensure robust watermark extraction even when the rendered samples are severly distorted (\eg, blurring, noise, and rotation). A random sampling strategy is further considered to make the protected model robust to different sampling strategy during rendering. 

Distortion-resistant rendering is only needed during the optimization of core models. If the core model is stolen, even with different rendering schemes and sampling strategies, the copyright message can still be robustly extracted. Our contribution can be summarized as follows:


%Instead of directly employing secret messages as part of the input, we propose to build an implict message representation for rendering. As displayed in~\Fref{fig:framework}, our message representation is able to \renjie{Something here}

%We strength the robustness of watermark extraction by considering the distortions in both 2D and 3D.   

%To guarantee the robustness when the rendered samples encounter external 2D distortions, we design a distortion-resistant rendering scheme for model training. By

%consider various 2D distortions during training. Then, 

%We propose a method


%In response to above issues, we propose a method (referred to as CopyRNeRF) to build an implicit message representation.  As outlined in~\Fref{fig:framework}, instead of directly employing secret messages as a part of the input, we select color, a component used for rendering in NeRF, as the cover medium~\cite{artz2001digital,qiu2022high} for the fusion of secret messages based on an implicit message representation. By integrating the spatial information, color, and messages as a single entity, the implicit message representation can successfully transfer the secret message into every rendered results while keeping their \textit{invisibility}. By considering spatial coordinates as a part of the representation, it also ensures the robust message extraction across different views. At last, during training, a distortion-resistant rendering is accommodated to better consider distortions from a wider patch scale, which ensures the messages can be \textit{robustly} extracted during testing. In summary, our major contribution can be concluded as follows:
\begin{itemize}
    \item a method to produce copyright-embedded NeRF models.
    
    \item a watermarked color representation to ensure invisibility and high rendering quality.
    
    %fusion strategy to invisibly fuse secret messages with the color selected as the cover medium.
    
    \item distortion-resistant rendering to ensure robustness across different rendering strategies or 2D distortions.
\end{itemize}
%\renjie{The remaining problem is how to }

%In order to tackle the problems mentioned above, we propose a novel two-stage NeRF framework and a patch-based training scheme.  An ordinary NeRF network is first trained to model the opacity field of the scene. Then the messages are only embedded to color field with density prediction fixed. There are two main reasons to do so. On the one hand, The opacity field of NeRF has the inherence of geometric consistency and can estimate the 3D coordinates of the rendered pixels, which can guarantee the strict consistency of the watermarked scene. On the other hand, messages can be hided in color channels in the rendering 2D images. Therefore, the geometry of 3D models will not have a big impact on the capacity of the watermarking. Besides, a patch-based training is introduced to the framework, aiming to enable the message extraction module to obtain information from an area of the image. This scheme can also improve the robustness of the watermarking, as the rendered images can be cropped into patches because of the distortion during the spread.

%Please follow the steps outlined below when submitting your manuscript to the IEEE Computer Society Press.
%This style guide now has several important modifications (for example, you are no longer warned against the use of sticky tape to attach your artwork to the paper), so all authors should read this new version.(a) Our method first generates the opacity $\sigma$ and the cover medium $\mathbf{c}$ by a \textit{medium generation} module. We input the concatenated color, spatial position, and viewing direction to build a color feature field. A message feature field is obtained by a network, and then fused with color feature field to get the message representation.
\begin{figure*}
  \centering
  \includegraphics[width=\linewidth]{supp_fig/framework_iccv_1.pdf}
  \caption{Illustration of our proposed method.  (a) A \textit{watermarked color representation} is obtained with the given secret message, which is able to produce watermarked color for rendering. (b) During training, a \textit{distortion-resistant rendering} is deployed to map the geometry ($\sigma$) and watermarked color representations to image patches with several distortions. (c) Finally, the secret message can be revealed by a CNN-based \textit{message extractor}. }
  
  %During sharing, we only need to share medium generation module and message representation module. The users can use volume rendering of NeRF to render 2D images with novel views. The message can still be robustly extracted.

  \label{fig:framework}
    \vspace{-10pt}
\end{figure*}

\section{Related work}

\paragraph{Neural radiance fields.} Various neural implicit scene representation schemes have been introduced recently \cite{niemeyer2020differentiable, yariv2020multiview, zhu2022neural}. The Scene Representation Networks (SNR) \cite{sitzmann2019scene} represent scenes as a multi-layer perceptron (MLP) that maps world coordinates to local features, which can be trained from 2D images and their camera poses. DeepSDF \cite{park2019deepsdf} and DIST \cite{liu2020dist} use trained networks to represent a continuous signed distance function of a class of shapes. PIFu \cite{saito2019pifu} learned two pixel-aligned implicit functions to infer surface and texture of clothed humans respectively from a single input image. Occupancy Networks \cite{mescheder2019occupancy, peng2020convolutional} are proposed as an implicit representation of 3D geometry of 3D objects or scenes with 3D supervision. 
NeRF~\cite{nerf2020, zhu2023occlusion} in particular directly maps the 3D position and 2D viewing direction to color and geometry by a MLP and synthesize novel views via volume rendering. The improvements and applications of this implicit representation have been rapidly growing in recent years, including NeRF accelerating~\cite{fridovich2022plenoxels, neff2021donerf}, sparse reconstruction~\cite{yu2021pixelnerf, chen2021mvsnerf}, and generative models~\cite{schwarz2020graf, chan2021pi}. NeRF models are not easy to train and may use private data, so protecting their copyright becomes crucial.


% There has been much promising recent work on using neural implicit scene representations. As these are continuous representations, they can easily handle complicated
% topologies and do not suffer from discretisation error, with
% the actual representative resolution depending on the capacity of the neural network used. The Scene Representation
% Network (SRN) [25] was one of the first methods to use a
% multi-layer perceptron (MLP) as the neural representation
% of a learned scene given a collection of images and associated poses. DeepSDF [19] and DIST [10] used deep decoders to learn implicit signed distance functions of various shape instances of the same class, and Occupancy Networks [15, 21] learned an implicit 3D occupancy function for shapes or large scale scenes given 3D supervision. NeRF [16] and other systems based on it [34, 11, 31, 27] use MLPs to overfit input from a single bounded scene
% and act as an implicit volumetric representation for realistic
% view-synthesis. 

\paragraph{Digital watermarking for 2D.} Early 2D watermarking approaches encode information in the least significant bits of image pixels \cite{413536}. Some other methods instead encode information in the transform domains \cite{lai2010digital}. Deep learning based methods for image watermarking have made substantial progress. HiDDeN~\cite{zhu2018hidden} was one of the first deep image watermarking methods that achieved superior performance compared to traditional watermarking approaches. RedMark~\cite{ahmadi2020redmark} introduced residual connections with a strength factor for
embedding binary images in the transform domain. Deep watermarking has since been generalized to video~\cite{weng2019high,zhang2019robust} as well. Modeling more complex and realistic image distortions also broadened the scope in terms of application~\cite{wengrowski2019light,tancik2020stegastamp}. However, those methods all cannot protect the copyright of 3D models.

\paragraph{Digital watermarking for 3D.} Traditional 3D watermarking approaches~\cite{ohbuchi2002frequency,praun1999robust,wu20153d} leveraged Fourier or wavelet analysis on triangular or polygonal meshes. Recently, Hou~\etal~\cite{hou2017blind} introduced a 3D watermarking method using the layering artifacts in 3D printed objects. Son~\etal~\cite{son2017perceptual} used mesh saliency as a perceptual metric to minimize vertex distortions. Hamidi~\etal~\cite{hamidi2019blind} further extended mesh saliency with wavelet transform to make 3D watermarking robust. Jing~\etal~\cite{liu2019novel} studied watermarking for point clouds through analyzing vertex curvatures. Recently, a deep-learning based approach \cite{Yoo_2022_CVPR} successfully embeds messages in 3D meshes and extracts them from 2D renderings. However, existing methods are for explicit 3D models, which cannot be used for NeRF models with implicit property.

\section{Preliminaries}
\label{sec:preliminary}
NeRF \cite{nerf2020} uses MLPs $\Theta_{\sigma}$ and $\Theta_{c}$ to map the 3D location $\mathbf{x} \in \mathbb{R}^3$ and viewing direction $\mathbf{d} \in \mathbb{R}^2$ to a color value $\mathbf{c} \in \mathbb{R}^3$ and a geometric value $\sigma \in \mathbb{R}^+$:
\begin{equation}
[\sigma, \mathbf{z}]=\Theta_{\sigma}\left(\gamma_{\mathbf{x}}(\mathbf{x})\right),
\label{eq:sigma}
\end{equation}
\begin{equation}
\mathbf{c}=\Theta_{c}\left(\mathbf{z}, \gamma_{\mathbf{d}}(\mathbf{d})\right),
\label{eq:nerffirst}
\end{equation}
where $\gamma_{\mathbf{x}}$ and $\gamma_{\mathbf{d}}$ are fixed encoding functions for location and viewing direction respectively. The intermediate variable $\mathbf{z}$ is a feature output by the first MLP $\Theta_{\sigma}$.

% To facilitate the fitting capability of the model, instead of directly using $\mathbf{x}$ and $\mathbf{d}$ as inputs to the MLP $F_\theta$, NeRF applies fixed positional encoding $\gamma(\cdot)$ to map the inputs to their Fourier features containing signals of muti-scale frequencies $\gamma(\mathbf{x})=\left[\sin{\mathbf{x}}, \cos{\mathbf{x}}, \cdots, \sin{2^{\Omega-1}\mathbf{x}}, \cos{2^{\Omega-1}\mathbf{x}} \right]^T$, where $\Omega$ is a hyper parameter controlling the bandwidth.

%
%\begin{equation}
  %\gamma(\mathbf{x})=\left[\sin{\mathbf{x}}, \cos{\mathbf{x}}, \cdots, \sin{2^{\Omega-1}\mathbf{x}}, \cos{2^{\Omega-1}\mathbf{x}} \right]^T,
%\label{fourier}
%\end{equation}
%

For rendering a 2D image from the radiance fields $\Theta_{\sigma}$ and $\Theta_{c}$, a numerical quadrature is used to approximate the volumetric projection integral. Formally, $N_p$ points are sampled along a camera ray $r$ with color and geometry values $\{(\mathbf{c}_r^i, \sigma_r^i)\}_{i=1}^N$. The RGB color value $\hat{\mathbf{C}}(r)$ is obtained using alpha composition
\begin{equation}
  \hat{\mathbf{C}}(r)=\sum_{i=1}^{N_p} T_r^i (1-\exp \left(-\sigma_r^i \delta_r^i\right)) \mathbf{c}_r^i,
  \label{eq:integral}
\end{equation}
where $T_r^i=\prod_{j=1}^{i-1}\left(\exp \left(-\sigma_r^i \delta_r^i\right)\right)$, and 
$\delta_r^i$ is the distance between adjacent sample points. The MLPs $\Theta_{\sigma}$ and $\Theta_{c}$ are optimized by minimizing a reconstruction loss between observations $\mathbf{C}$ and predictions $\hat{\mathbf{C}}$ as
\begin{equation}
  \mathcal{L}_{recon} = \frac{1}{N_r}\sum_{m =1}^{N_r}\|\hat{\mathbf{C}}(r_m)-\mathbf{C}(r_m)\|_2^2,
  \label{eq:recon}
\end{equation}
where $N_r$ is the number of sampling pixels. Given $\Theta_{\sigma}$ and $\Theta_{c}$, novel views can be synthesized by invoking volume rendering for each ray.


%NeRF provides a convenient way to obtain an implicit 3D representation via learning from a collection of 2D images. 
Considering the superior capability of NeRF in rendering novel views and representing various scenes, how can we protect its copyright when it is stolen by malicious users?

%and identify the ownership of any rendered 2D images using this representation?

%This representation can be rendered from any viewpoint via differentiable volume
%rendering, and hence can be optimized to fit a collection of input photos captured
%from multiple views, and then later used for synthesizing photo-realistic novel
%views. We move beyond photo-realism and add watermark to the radiance
%field by embedding it with messages for copyright protection.
% The rendered pixel value for camera ray $r$ can then be compared against the corresponding ground truth pixel value, $C(r)$, for all the camera rays of the target view with pose $P$. The NeRF rendering loss is thus given by
% \begin{equation}
%   L = \sum_{r \in R(P)}\|\hat{C}(r)-C(r)\|_2^2,
%   \label{eq:integral}
% \end{equation}
% where $R(P)$ is the set of all camera rays of target pose $P$. and its corresponding camera parameters $\{(R_n,t_n)\}_{n=1}^N$Then, the secret messages implicitly embedded into the NeRF representation can be rendered out along with the novel view synthesis, which ensures the copyright protection of this implicit 3D model. As outlined in Fig.XXX, with a collection of 2D images $\{I_n\}_{n=1}^N$ and the binary message $M \in \{0,1\}^{N_b}$ with length $N_b$, our goal is to embed the message $M$ into the NeRF model and extract it from the rendered 2D images in any views to verify its copyright.We address the issue raised in~\Sref{sec:preliminary} by ensuring every 2D images rendered from such an implicit representation are watermarked while having high quality. 

\section{Proposed method}

%f
As outlined in~\Fref{fig:framework}, with a collection of 2D images $\{I_n\}_{n=1}^N$ and the binary message $\M \in \{0,1\}^{N_b}$ with length $N_b$, we address the issue raised in~\Sref{sec:preliminary} by building a watermarked color representation during optimization. In training, a distortion-resistant rendering is further applied to improve the robustness when 2D distortions or different rendering schemes are encountered. With the above design, the secret messages can be robustly extracted during testing even 
encountering sever distortions or different rendering strategies.
%using a conventional CNN-based message extractor or a different rendering scheme. 

%our goal is to embed the message $M$ into the NeRF model learned from the 2D image collection, and further extract the secret messages from the rendered 2D images in any views to verify its copyright.

%We illustrate our framework for watermarking a neural implicit representation from a collection of images $\{I_n\}_{n=1}^N$ with corresponding camera parameters $\{(R_n,t_n)\}_{n=1}^N$. Given the binary message $M \in \{0,1\}^{N_b}$ with length $N_b$, our goal is to embed the message $M$ into NeRF and extract it from the rendered 2D images in any views while keeping high-quality reconstruction of 3D content. As we discussed in \cref{sec:intro}, although both opacity and color in implicit 3D representation can be embedded with messages, we only embed messages in the radiance color field, as geometry information is difficult
%to be directly accessed from 2D images. Besides, messages cannot be effectively accessed by an extractor designed for the whole image if pixel-based rendering is deployed. We propose a novel robust watermarking approach for 3D scenes. In the next, we will first introduce the watermarking implicit 3D representation in \cref{}, and then discuss our distortion robust rendering in \cref{}. After that, the message extractor is introduced in \cref{}. Finally, we describe how our model is trained in \cref{}.Several methods have been proposed to embed information into NeRF-based representation. The common way is proposed by PixelNeRF~\cite{yu2021pixelnerf}, which directly concatenate the additional features with the vanilla coordinates used as the input. However, such a way may significantly influence the appearance of rendered novel views, which makes observers easily feel the existence of additional information. This can avoid the ambiguity caused by the geometry~\cite{bas2019does} during message extraction in 2D.


\subsection{Building watermarked color representation}
\label{sec:4.1}
The rendering in \Eref{eq:integral} relies on color and geometry produced by their corresponding representation in NeRF. To ensure the transmission of copyright messages to the rendered results, we propose embedding messages into their representation. We create a watermarked color representation on the basis of $\Theta_c$ defined in~\Eref{eq:nerffirst} to guarantee the message invisibility and consistency across viewpoints. The representation of geometry is also the potential for watermarking, but external information on geometry may undermine rendering quality~\cite{wang2022nerf, Huang22StylizedNeRF, chiang2022stylizing}. Therefore, the geometry does not become our first option, while experiments are also conducted to verify this setting.

We keep the geometry representation in~\Eref{eq:sigma} unchanged, and construct the watermarked color representation $\Theta_{m}$ to produce the message embedded color $\mathbf{c}_m$ as follows:
%\begin{equation}
%[\sigma, \mathbf{z}]=\operatorname{MLP}_{\theta_1}\left(\gamma_{\mathbf{x}}(\mathbf{r}(t))\right)
%\end{equation}
\begin{equation}
\mathbf{c}_m=\Theta_{m}\left(\mathbf{c},\gamma_{\mathbf{x}}(\mathbf{x}), \gamma_{\mathbf{d}}(\mathbf{d}), \mathbf{M}\right),
\end{equation}
where $\mathbf{M}$ denotes the message to be embedded and $\Theta_{m}$ contains several MLPs to ensure reliable message embedding. The input $\mathbf{c}$ is obtained by querying $\Theta_{c}$ using~\Eref{eq:nerffirst}. Several previous methods have pointed out the importance of building a 3D feature field when distributed features are needed to characterize 
composite information~\cite{xu20223d,chan2022efficient}. Thus, instead of directly fusing those information, we first construct their corresponding feature field and then combine them progressively.

%Specifically, to ensure the robust message extraction across different viewpoints, we still consider the 3D localtion $\mathbf{x}$ as a part of the input for watermarked color representation, instead of solely employing viewpoints $\mathbf{d}$ as~\Eref{eq:nerffirst}.

%experiments are also conducted to verify this message in geometry setting. 

%Thus, the geometry does not become our first option, while experiments are also conducted to verify this message in geometry setting. 
%exponential operation applied on geometry in~\Eref{eq:integral} may make the message extraction difficult 

%The volume rendering used by NeRF relies on the color and geometry produced by their corresponding represenation. We  

%Our message representation is designed to directly produces factors enough for subsequent rendering

%To guarantee invisibility within rendered results, we watermark a subset of representation as the message representation for subsequent rendering.

%a lighweigt CNN-based message extractor can be leveraged to access secret messages from rendering samples. 

%The message repres

%Then, the messages representation can directly produce 

%Our message fusion module aims at creating an implicit message representation for the subsequent rendering. Then, even a lightweight CNN-based message extractor can be leveraged to access secret messages from synthesized examples during testing. In contrast to the feature embedding scheme in~\cite{yu2021pixelnerf} that directly manipulates the implicit representation related to scene details, we select color as the cover medium to build an implicit message representation. By fixing geometry, our method can guarantee higher visual quality in the synthesized results. 

%By fusing secret messages, color, and spatial position into a single entity, our message representation ensures the invisibility in the rendered results. 

%With the accommodated distortion-resistant rendering, our method can further guarantee the robust extraction under different distortions even with a simple CNN-based message extractor.


% \paragraph{Color embedding scheme}

% We strictly follow the the established requirement for invisibility and robustness in digital watermarking when designing our framework. Though several methods have been proposed to embed information into NeRF-based representation, one key issue is still to be addressed: How to preserve the geometry consistency during the message embedding. As a message .

% they still face one challenge about how to preserve the geometry consistency of the rendered results. As xxxx, 


% A color embedding mechanism is introduced to implicitly embed the secret messages into the network parameters. Then, within this scheme,  by indirectly affecting the color component, our proposed color embedding scheme can both effectively embed messages into the implicit network parameters and ensure its invisibility at its rendered 2D counterparts. First, one key issue is how to ensure the invisibility of the embedded secret messages. Though 



%watermarking framework for NeRF strictly follows the established requirement for invisibility and robustness in digital watermarking. An implicit watermark embedding scheme is introduced in~\cref{sec:4.1}, where we embed the messages into the implicit 3D representation by manipulating the color component used for volume rendering. Then, a distortion-robust rendering framework is proposed in~\cref{sec:4.2}. Despite the invisible embedding, the digital watermarking should also be effectively extracted even the render 2D image is distorted. Most distortion in 2D is a global operation beyond the capability of pixel-based mapping in NeRF. We introduce a distortion-robust rendering framework by aggregating information from a wider area, which guarantee the robust message extraction when encountering such global distortion.


%\renjie{Our implicit watermark embedding scheme aims at hiding the message into the NeRF-based 3D representation. To address the requirement for invisibility,   }

%We use a MLP to generate a higher-dimensional 3D feature volume to better describe the complex 3D scenes (other reasons). Besides, a preparation network is adopted to transform secret messages into features. The spatial feature filed is then concatenated with message features to produce the watermarked color. Besides, a skip connection is added to further guarantee the model's stability.
%We propose a watermarking framework to embed messages in radiance color field of the implicit representation with the opacity field fixed. Therefore, the implicit 3D representation with watermark can be obtained in two steps. The opacity field of the target scene, as well as the original color field, are firstly generated by a NeRF model from a set of 2D images. Then the message $M$ is embedded into color field by a encoder block while keeping the geometry fixed.Our target is to embed secret messages into the NeRF-based implicit 3D representation. Then, we can easily extract the embedded messages from the rendered 2D images by using even a CNN-based message extractor with the simple structure. To preserve the geometry of the rendered result, we first use a geometry generator based on the ordinary NeRF to learn the intermediate factor necessary for message fusion, and then embed the messages into color based on a message representation.

%\paragraph{Base representation}\label{sec:geo}
%We first employ a medium generation module to generate the cover mediums necessary for message fusion. Since we select color defined in~\Eref{eq:nerffirst} as the medium, the whole generation process follows the setting in the ordinary NeRF~\cite{nerf2020}. We randomly sample $N_r$ pixel from input images to form a batch of camera rays and 3D points marching along the rays based on the corresponding camera poses and camera intrinsic. Then, by using a MLP, we can estimate the color value $\mathbf{c}$ and geometry value $\sigma$. Then, a standard volume rendering approach is applied to render the RGB color value $\hat{\mathbf{C}}(r_m)$ of each ray $r_m$. We can further optimize the two factors by training MLPs using a mean square error between the RGB color value $\hat{\mathbf{C}}(r_m)$ and its corresponding ground truth color $\mathbf{C}(r_m)$ of the corresponding image pixel of ray $r_m$ as  
%\begin{equation}
%  \mathcal{L}_m = \frac{1}{N_r}\sum_{m =1}^{N_r}\|\hat{\mathbf{C}}(r_m)-\mathbf{C}(r_m)\|_2^2.
%\end{equation}
%The objective for training the MLPs of the scene function is the mean square error between $\hat{C}(r_m)$ and its ground truth $C(r_m)$ of the corresponding image pixel of ray $r_m$: 
%\begin{equation}
%  \mathcal{L}_g = \frac{1}{N_r}\sum_{m =1}^{N_r}\|\hat{C}(r_m)-C(r_m)\|_2^2.
%\end{equation}
%Though selecting $\sigma$ as the cover medium is also an option, the geometry in 2D views is with ambiguity~\cite{bas2019does}, which may undermine the message extraction. Therefore, we fix the geometry and focus on fusing the message into the color $\mathbf{c}$ by using an implicit message representation. 

%Although NeRF \cite{10.1007/978-3-030-58452-8_24} model is adopted in this stage, it should be noted that any versions of NeRF structure can be deployed here. The training process is briefly summarized as follows. At each optimization iteration, we randomly sample $M$ pixels from input images to form a batch of camera rays $\{r_m\}^M_{m=1}$ and 3D points marching along the rays based on the corresponding camera poses and camera intrinsic. With taking these 3D points and their corresponding view directions, \ie the directions of the corresponding camera rays, as input to the scene function of NeRF to obtain the output set of volume opacities and colors, the volume rendering technique is used to render the color value $\hat{C}(r_m)$ of each ray $r_m$. The objective for training the MLPs of the scene function is the mean square error (MSE) between $\hat{C}(r_m)$ and the groundtruth color $C(r_m)$ of the corresponding image pixel of ray $r_m$. 
%(copied from Stylizing 3D Scene via Implicit Representation and HyperNetwork)
%\begin{equation}
%  L_1 = \frac{1}{M}\sum_{m =1}^M\|\hat{C}(r_m)-C(r_m)\|_2^2
%\end{equation}


%Therefore, nothing is more substantial that building a representation in advance. 

%In this module, we aim to learn a neural representation of the target 3D scene from the given images taken from different camera poses. Although NeRF \cite{10.1007/978-3-030-58452-8_24} model is adopted in this stage, it should be noted that any types of NeRF structure can be deployed here. When querying the network at point $(x,d)$, the opacity $\sigma(x)$ and the original color $c_o(x,d)$ can be obtained. 

% The training process is briefly summarized as follows. At each optimization iteration, we randomly sample $M$ pixels from input images to form a batch of camera rays $\{r_m\}^M_{m=1}$ and 3D points marching along the rays based on the corresponding camera poses and camera intrinsic. With taking these 3D points and their corresponding view directions, \ie the directions of the corresponding camera rays, as input to the scene function of NeRF to obtain the output set of volume opacities and colors, the volume rendering technique is used to render the color value $\hat{C}(r_m)$ of each ray $r_m$. The objective for training the MLPs of the scene function is the mean square error (MSE) between $\hat{C}(r_m)$ and the groundtruth color $C(r_m)$ of the corresponding image pixel of ray $r_m$:
% %(copied from Stylizing 3D Scene via Implicit Representation and HyperNetwork)
% \begin{equation}
%   L_1 = \frac{1}{M}\sum_{m =1}^M\|\hat{C}(r_m)-C(r_m)\|_2^2
%   \label{eq:mse}
% \end{equation}

%\subsubsection{Building implicit message representation}
%As discussed in previous sections, despite directly manipulating the network parameters using secret messages, we focus on embedding message into the cover medium by using an implicit message representation constructed by fusing color feature field and message feature field. 

%NeRF is highly sensitive to the new information injected into it. To preserve the invisibility of the embedded messages, we fully control the geometry component and focus on how to inject the messages into the color component.The implicit 3D scene representation obtained in~\Sref{sec:geo} has encoded the complete geometry and the original appearance of the target scenes into a corresponding NeRF model.

%As in the previous geometry generation step, we have encoded the complete geometry and the original appearance of the target 3D scene into the NeRF model. We focus on learning the message embedding module to further predict the watermarked colors with the given message $M$. As we only embed messages in the color component, the geometry generation module and the original color generation module are fixed during the secret message embedding. This embedding process is designed to be implemented by three modules, spatial feature extractor $E_{\theta}$, message feature extractor $F_{\phi}$, and message-embedded color generator $G_{\lambda}$, where $\theta$, $\phi$, and $\lambda$ are trainable parameters. By using appropriate training schemes, directly combining messages with color can also make the messages invisible at the rendered results.

\noindent\textbf{Color feature field.} In this stage, we aim at fusing the spatial information and color representation to ensure message consistency and robustness across viewpoints.  We adopt a color feature field by considering color, spatial positions, and viewing directions simultaneously as follows:
\begin{equation}
f_c=E_\xi(\mathbf{c},\gamma_{\mathbf{x}}(\mathbf{x}),\gamma_{\mathbf{d}}(\mathbf{d})).
% \begin{aligned}
% E_\theta: &\mathbb{R}^{3} \times \mathbb{R}^{\Omega_{\mathbf{x}}} \times \mathbb{R}^{\Omega_{\mathbf{d}}} \rightarrow \mathbb{R}^{N_s},\\
% &(\mathbf{c}, \gamma(), \gamma(\mathbf{d})) \mapsto .
% \end{aligned}
\label{eq:colorrepresentation}
\end{equation}
%
Given a 3D coordinate $\mathbf{x}$ and a viewing direction $\mathbf{d}$, we first query the color representation $\Theta_{c}\left(\mathbf{z}, \gamma_{\mathbf{d}}(\mathbf{d})\right)$ to get $\mathbf{c}$, and then concatenate them with $\mathbf{x}$ and $\mathbf{d}$ to obtain spatial descriptor $\mathbf{v}$ as the input. Then the color feature encoder $E_{\xi}$ transforms $\mathbf{v}$ to the high-dimensional color feature field $f_c$ with dimension $N_c$. The Fourier feature encoding is applied to $\mathbf{x}$ and $\mathbf{d}$ before the feature extraction.

%this can never ensure the robust messages 


%It has been proved that the digital watermarking can achieve better performance by extracting both high- and low-level features~\cite{9133124}. 
%We use a module to extract spatial feature rather than directly changing the color of the point. Such a spatial feature extractor can perform watermarking in a more adaptive manner by dynamically learning algorithms to extract both high-level and low-level features for watermarking~\cite{9133124}. Different from the CNN-based feature extracting, which is implemented in the image level, we generate a coordinate-based volume feature. Given a 3D coordinate $x$ and a viewing direction $d$, we first query the opacity $\sigma(x)$ and color $c_o(x,d)$ from the (opacity volume) and (basic color volume), and then concatenate them with $x$ and $d$ to obtain spatial descriptor $v$ as the input. Then the spatial feature extractor $E_{\theta}$ transforms $v$ to the high-dimensional spatial feature field $f_s$ with dimension $N_s$. The Fourier feature encoding $\gamma$ in~\Eref{fourier} is applied to $x$ and $d$ before feature extracting. The above process can be formulated as
%\begin{equation}
%v={\rm Concat} (c_o, \sigma, \gamma(x), \gamma(d)),
%\end{equation}
%\begin{equation}
%f_s(x,d)=E_{\theta}(v).
%\end{equation}

%or

%\begin{equation}
%\begin{aligned}
%E_\theta: &\mathbb{R}^{L_{\mathbf{x}}} \times \mathbb{R}^{+} \times %\mathbb{R}^{L_{\mathbf{x}}} \times \mathbb{R}^{L_{\mathbf{d}}} %\rightarrow \mathbb{R}^{N_s},\\
%&(\mathbf{c}_o, \sigma, \gamma(\mathbf{x}), \gamma(\mathbf{d})) \mapsto %f_s,
%\end{aligned}
%\end{equation}
%where ${L_{\mathbf{x}}}$ and ${L_{\mathbf{d}}}$ is the dimension of Fourier features of $\mathbf{x}$ and $\mathbf{d}$ respectively.

\noindent\textbf{Message feature field.} We further construct the message feature field. Specifically, we follow the classical setting in digital watermarking by transforming secret messages into higher dimensions~\cite{NIPS2017_838e8afb, baluja2019hiding}. It ensures more succinctly encoding of desired messages~\cite{NIPS2017_838e8afb}. As shown in~\Fref{fig:framework}, a message feature encoder is applied to map the messages to its corresponding higher dimensions as follows:
\begin{equation}
f_\M=D_{\phi}(\M).
\label{eq:messageembedding}
\end{equation}
In~\Eref{eq:messageembedding}, given message $\M$ of length $N_b$, the message feature encoder $D_{\phi}$ applies a MLP to the input message, resulting in a message feature field $f_\M$ of dimension $N_{m}$. 

%The message feature $f_m$ is then replicated to form a message feature field $f_\M$ in 3D space.



Then, the watermarked color can be generated via a feature fusion module $G_{\psi}$ that integrates both color feature field and message feature field as follows: 
\begin{equation}
\mathbf{c}_m=G_{\psi}(f_c,f_\M,\mathbf{c}).
\label{eq:messagerepresentation}
\end{equation}
Specifically, $\mathbf{c}$ is also employed here to make the final results more stable. $\mathbf{c}_m$ is with the same dimension to $\mathbf{c}$, which ensures this representation can easily adapt to current rendering schemes.


%This structure has several nice properties for implicit watermarking in this works. First, it grants higher flexibility to our framework as it can be applied to variour variants

%We adopt the idea of preparation networks in 2D image~\cite{NIPS2017_838e8afb}, which is to transform the secret information to more useful features for succinctly encoding. Here a message feature extractor is applied to map the message to the feature in a higher dimensional space. Given message $M$ of length $N_b$, the message feature extractor $F_{\phi}$ applies a MLP to the input message, resulting in a feature of dimension $N_{m}$.
%\begin{equation}
%f_m(x,d)=F_{\phi}(M).
%\end{equation}


%\paragraph{Neural message fusion}
%We use a neural feature fusion block to combine the spatial feature field $f_s$ and the message feature field $f_m$, and generate a new color volume field $c_m$, which has message embedded in. Besides these two features, the basic color volume $c_o$ generated in~\Sref{geo} is also used as input of the networks, playing the role of skip connection. As described in~\ref{}, the skip connection can make the networks architecture more stable. Therefore, the color volume filed can be denoted as
%\begin{equation}
%c_m(x,d)=G_{\lambda}(f_s,f_m,c_o).
%\end{equation}


% As illustrated in Fig. XXX, the message embedding block takes the original color $c_o(x,d)$, message $M$, opacity $\sigma(x)$ and $(x,d)$ as input to output the watermarked color $c_m(c_o,M,\sigma,x,d)$.

%Unlike the methods of style transfer NeRF that modify the color branch of the original networks directly, here an additional message embedding block is added. Such structure benefits the watermarking task in two reasons. Firstly, this structure has higher flexibility as it can be applied to various of 3D neural representations. Secondly, this structure enable us to use the skip connection that connects the original color $c_o(x,d)$ to the output layer of the networks, which makes the overall architecture more stable.


\subsection{Distortion-resistant rendering}
\label{sec:4.2}
Directly employing the watermarked representation for volume rendering has already been able to guarantee invisibility and robustness across viewpoints. However, as discussed in~\Sref{sec:intro}, the message should be robustly extracted even when encountering diverse distortion to the rendered 2D images. Besides, for an implicit model relying on rendering to display its contents, the robustness should also be secured even when different rendering strategies are employed. Such requirement for robustness cannot be achieved by simply using  watermarked representation under the classical NeRF training framework. For example, the pixel-wise rendering strategy cannot effectively model the distortion (\eg, blurring and cropping) only meaningful in a wider scale. We, therefore, propose a distortion-resistant rendering by strengthing the robustness using a random sampling strategy and distortion layer.

%the digital watermarking should also be effectively extracted even the rendered 2D images are distorted. Despite the effectiveness of NeRF in learning 3D scene representation from 2D images, it actually builds a local mapping correlation by mapping a spatial position to its corresponding image pixels. However, most distortion in 2D is a global operation XX. To enable encoded messages can survive distortions of various forms, we propose a distortion robust rendering to improve the robustness of the model. extend the idea of patching from traditional 3D watermarking methods~\cite{kwon2003watermarking, cotting2004robust} to the volume rendering process.

Since most 2D distortions can only be obviously observed in a certain area, we consider the rendering process in a patch level~\cite{kwon2003watermarking, cotting2004robust}. A window with the random position is cropped from the input image with a certain height and width, then we uniformly sample the pixels from such window to form a smaller patch. The center of the patch is denoted by $\mathbf{u}=(u, v) \in \mathbb{R}^2$, and the size of patch is determined by $K\in \mathbb{R}^{+}$. We randomly draw the patch center $\mathbf{u}$ from a uniform distribution $\mathbf{u} \sim \mathcal{U}(\Omega)$ over the image domain $\Omega$. The patch $\mathcal{P}(\mathbf{u}, K)$ can be denoted by by a set of 2D image coordinates as
\begin{equation}
\mathcal{P}(\mathbf{u}, K)=\{(x+u, y+v) \mid x, y \in\{-\frac{K}{2}, \ldots, \frac{K}{2}-1\}\}.
\end{equation}
Such a patch-based scheme constitutes the backbone of our distortion-resistant rendering, due to its advantages in capturing information on a wider scale. Specifically, we employ a variable patch size to accommodate diverse distortions during rendering, which can ensure higher robustness in message extraction. This is because small patches increase the robustness against cropping attacks and large patches allow higher redundancy in the bit encoding, which leads to increased resilience against random noise~\cite{cotting2004robust}.




%Partitioning benefits the watermarking task for several reasons.

%(1) From a rendering perspective, messages are supposed to be hidden in an area but not in a single pixel of the rendered 2D image. Messages can therefore be extracted by a CNN-based network.

%(2) Using multiple patches allows a watermark to be extracted even in case of large deformations and significant attacks on the signed object as long as one single patch is sufficiently preserved and can be registered with the original object. Patching also helps when watermarks have to be extracted from cropped objects.

%(3) A variable patch size can accommodate different
%security requirements: On the one hand, small patches
%increase the robustness against cropping attacks. On
%the other hand, large patches allow higher redundancy
%in the bit encoding, leading to increased resilience
%against random noise. A variable, user-defined patch
%size is therefore desirable.

%(4) Patches ensure the scalability of the watermarking procedure. The time and memory efficiency of the algorithm using patches is higher than for the same procedure defined over the entire image. This is due to the fact that NeRF needs to query hundreds of sample points along the ray to render a single pixel.

As the corresponding 3D rays are uniquely determined by $\mathcal{P}(\mathbf{u}, K)$, the camera pose and intrinsics, the image patch $\widetilde{\mathbf{P}}$ can be obtained after points sampling and rendering. Based on the sampling points in~\Sref{sec:preliminary}, we use a random sampling scheme to further improve the model's robustness, which is described as follows.

\noindent\textbf{Random sampling.} During volume rendering, NeRF~\cite{nerf2020} is required to sample 3D points along a ray to calculate the RGB value of a pixel color. However, the sampling strategy may vary as the renderer changes~\cite{neff2021donerf, lindell2021autoint}. To make our message extraction more robust even under different sampling strategies, we employ a random sampling strategy by adding a shifting value to the sampling points. Specifically, the original $N_p$ sampling points along ray $r$ is denoted by a sequence, which can be concluded as $\mathcal{X} = (x_r^1, x_r^2, \cdots, x_r^{N_p})$, where $x_r^{i},  i = 1, 2, \cdots, N_p$ denotes the sampling points during rendering. The randomized sample sequence $\mathcal{X}_{random}$ can be denoted by adding a shifting value as
\begin{equation}
\begin{aligned}
\mathcal{X}_{random} = (x_r^1 + z^1, x_r^2 + z^2, \cdots, x_r^{N_p} + z^{N_p}),\\
z^i \sim \mathcal{N}(0, \beta^2),\ i=1,2, \cdots, N_p,
\end{aligned}
\end{equation}
where $\mathcal{N}(0, \beta^2)$ is the Gaussian distribution with zero mean and standard deviation $\beta$.

By querying the watermarked color representation and geometry values at  $N_p$ points in $\mathcal{X}_{random}$, the rendering operator can be then applied to generate the watermarked color $\widetilde{\mathbf{C}}_m$ in rendered images:
\begin{equation}
  \widetilde{\mathbf{C}}_m(r)=\sum_{i=1}^{N_p} T_r^i (1-\exp \left(-\sigma_r^i \delta_r^i\right)) \mathbf{c}_m^i,
\end{equation}
where $T_r^i$ and $\delta_r^i$ are with the same definitions to their counterparts in~\Eref{eq:integral}.
%=\left\|(\mathbf{x}_r^{i+1} + z^{i+1})-(\mathbf{x}_r^i+z^{i+1})\right\|_2

All the colors obtained by  coordinates $\mathcal{P}$ can form a $K\times K$ image patch $\widetilde{\mathbf{P}}$. The content loss $\mathcal{L}_{content}$ of the 3D representation is calculated between watermarked patch $\widetilde{\mathbf{P}}$ and the $\hat{\mathbf{P}}$, where $\hat{\mathbf{P}}$ is rendered from the non-watermarked representation by the same coordinates $\mathcal{P}$.
In detail, the content loss $\mathcal{L}_{content}$ has two components namely pixel-wise MSE loss and perceptual loss:
\begin{equation}
  \mathcal{L}_{content} = \|\widetilde{\mathbf{P}}-\hat{\mathbf{P}}\|_2^2 + \lambda\|\Psi(\widetilde{\mathbf{P}}) - \Psi(\hat{\mathbf{P}})\|_2^2,
  \label{eq:content}
\end{equation}
where $\Psi(\cdot)$ denotes the feature representation obtained from a VGG-16 network, and $\lambda$ is a hyperparameter to balance the loss functions. 

\noindent\textbf{Distortion layer.}
To make our watermarking system robust to 2D distortions, a distortion layer is employed in our watermarking training pipeline after the patch $\widetilde{\mathbf{P}}$ is rendered. Several commonly used distortions are considered: 1) additive Gaussian noise with mean $\mu$ and standard deviation $\nu$; 2) random axis-angle rotation with parameters $\alpha$; and 3) random scaling with a parameter $s$; 4) Gaussian blur with kernel $k$. Since all these distortions are differentiable, we could train our network end-to-end. 

The distortion-resistant rendering is only applied during training. It is not a part of the core model. If the core model is stolen, even malicous users use different rendering strategy, the expected robustness can still be secured.

%Once shared the  with others, the users can also employ classical volume rendering in vanilla NeRF~\cite{nerf2020} to get the results, and our message can also be robustly extracted. 

% More formally, we have:
% \begin{align}
% &\widetilde{\mathbf{P}}_{noise} = \widetilde{\mathbf{P}} + n,\ n \sim \mathcal{N}(\mu, \nu^2),\\
% &\widetilde{\mathbf{P}}_{rotation} = \widetilde{\mathbf{P}} \cdot \mathbf{R}_{\alpha},\\
% &\widetilde{\mathbf{P}}_{scale} = {\rm Resize}(\widetilde{\mathbf{P}}, s),\\
% &\widetilde{\mathbf{P}}_{blur} = \widetilde{\mathbf{P}} \otimes k,
% \end{align}
% where $\mathbf{R}_{\alpha}$ is a rotation matrix, and $\otimes$ is the convolution operator.
% \subsection{Message Extraction}

\subsection{Message extractor}
To retrieve message $\hat{\M}$ from the $K \times K$ rendered patch $\mathbf{P}$, a message extractor $H_{\chi}$ is proposed to be trained end-to-end:
\begin{equation}
H_{\chi}: \mathbb{R}^{K \times K}\rightarrow \mathbb{R}^{N_b},\ \mathbf{P} \mapsto \hat{\M},
\end{equation}
where $\chi$ is a trainable parameter. 
Specifically, we employ a sequence of 2D convolutional layers with the batch normalization and ReLU functions~\cite{ioffe2015batch}. An average pooling is then performed, following by a final linear layer with a fixed output dimension $N_b$, which is the length of the message, to produce the continuous predicted message $\hat{\M}$. Because of the use of average pooling, the message extractor is compatible with any patch sizes, which means the network structure can remain unchanged when applying size-changing distortions such as random scaling.

The message loss $\mathcal{L}_m$ is then obtained by calculating the mean square error between predicted message $\hat{\mathbf{M}}$ and the ground truth message $\mathbf{M}$:
\begin{equation}
  \mathcal{L}_m = \|\hat{\mathbf{M}}-\mathbf{M}\|_2^2.
\end{equation}

To evaluate the bit accuracy during testing, the binary predicted message $\hat{\mathbf{M}}_b$ can be obtained by rounding:
\begin{equation}
\hat{\mathbf{M}}_b =  {\rm clamp}({\rm sign}(\hat{\mathbf{M}}), 0, 1),
\end{equation}
% The bit accuracy can be calculated as
% \begin{equation}
% \hat{\mathbf{M}}_b =  {\rm Clamp}({\rm Round}(\hat{\mathbf{M}}), 0, 1).
% \end{equation}
where $\rm clamp$ and $\rm sign$ are of the same definitions in~\cite{Yoo_2022_CVPR}. It should be noted that we use the continuous result $\hat{\mathbf{M}}$ in the training process, while the binary one $\hat{\mathbf{M}}_b$ is only adopted in testing process.

Therefore, the overall loss to train the copyright-protected neural radiance fields can be obtained as
\begin{equation}
\mathcal{L} = \gamma_1\mathcal{L}_{content} + \gamma_2\mathcal{L}_m,
\label{eq:loss}
\end{equation}
where $\gamma_1$ and $\gamma_2$ are hyperparameters to balance the loss functions.


\begin{figure*}
  \centering
  \includegraphics[width=\linewidth]{figs/fig_e1_5.pdf}
  \caption{Visual quality comparisons of each baseline. We show the differences ($\times10$) between the synthesized results and the ground truth next to each method. Our proposed \CopyNeRF can achieve a well balance between the reconstruction quality and bit accuracy.} %Hidden~\cite{zhu2018hidden} + NeRF~\cite{nerf2020} can achieve higher PSNR, while its bit accuracy is very low. NeRF with messages can achieve higher bit accuracy when compared with Hidden~\cite{zhu2018hidden} + NeRF~\cite{nerf2020}, while its embedding scheme deteriorate its PSNR values.
  \label{fig:visualquality}
  \vspace{-5pt}
\end{figure*}

% \subsection{Training Strategy}
% To better reserve the geometry of the 3D representations, 
% \paragraph{Geometric Training Stage (First Stage).}
% \begin{equation}
%   L_1 = \frac{1}{M}\sum_{m =1}^M\|\hat{C}(r_m)-C(r_m)\|_2^2
%   \label{eq:mse}
% \end{equation}

% \paragraph{Watermarking Training Stage (Second Stage).}
% Using our patch-based algorithm to select the patch $P(I_n)$
% from input image $I_n$ and its paired patch $P(\widetilde{I}_n)$ from the watermarked result $\widetilde{I}_n$, the content loss is computed between them. In detail, the content loss has two component. One is the pixel-wise MSE loss $L_{mse}$, which can be calculated as \Eref{eq:mse}. The other content is perceptual loss $L_{percep}$, which is defined as:
% \begin{equation}
%   L_{percep} = \| - \|_2^2,
%   \label{eq:percep}
% \end{equation}
% where XX denotes the feature representation obtained from the relu4\_1 layer of an ImageNet-pretrained VGG-16 network. 

% Message loss:


%\subsection{Geometry Generation.}
%In this module we aim to learn a NeRF-based representation of the target 3D scene from the given $N$ images $\{I_n\}^N_{n=1}$ taken from different camera poses $\{(R_n, t_n)\}^N_{n=1}$. For simplicity, we directly use the original NeRF as the background to learn the corresponding scene representation. However, based on our experiences, any NeRF variants can be considered in this place. 


%Although NeRF \cite{10.1007/978-3-030-58452-8_24} model is adopted in this stage, it should be noted that any versions of NeRF structure can be deployed here. The training process is briefly summarized as follows. At each optimization iteration, we randomly sample $M$ pixels from input images to form a batch of camera rays $\{r_m\}^M_{m=1}$ and 3D points marching along the rays based on the corresponding camera poses and camera intrinsic. With taking these 3D points and their corresponding view directions, \ie the directions of the corresponding camera rays, as input to the scene function of NeRF to obtain the output set of volume opacities and colors, the volume rendering technique is used to render the color value $\hat{C}(r_m)$ of each ray $r_m$. The objective for training the MLPs of the scene function is the mean square error (MSE) between $\hat{C}(r_m)$ and the groundtruth color $C(r_m)$ of the corresponding image pixel of ray $r_m$. 
%(copied from Stylizing 3D Scene via Implicit Representation and HyperNetwork)
%\begin{equation}
%  L_1 = \frac{1}{M}\sum_{m =1}^M\|\hat{C}(r_m)-C(r_m)\|_2^2
%\end{equation}

%\subsection{Message Embedding.}


%\Fref{} shows the network architecture of message embedding module. 

%\subsection{Patch-based rendering.}
%We extend the idea of patching from traditional 3D watermarking methods \cite{kwon2003watermarking, cotting2004robust} to the volume rendering process. A random window is randomly cropped from the input image with a certain height and width, then we uniformly sample the pixels from such window to form a smaller patch. 

%Partitioning benefits the watermarking task for several reasons.

%(1) From a rendering perspective, messages are supposed to be hidden in an area but not in a single pixel of the rendered 2D image. Messages can therefore be extracted by the CNN-based networks.

%(2) Using multiple patches allows a watermark to be extracted even in case of large deformations and significant attacks on the signed object as long as one single patch is sufficiently preserved and can be registered with the original object. Patching also helps when watermarks have to be extracted from cropped objects.

%(3) A variable patch size can accommodate different
%security requirements: On the one hand, small patches
%increase the robustness against cropping attacks. On
%the other hand, large patches allow higher redundancy
%in the bit encoding, leading to increased resilience
%size is therefore desirable.

%(4) Patches ensure the scalability of the watermarking procedure. The time and memory efficiency of the algorithm using patches is higher than for the same procedure defined over the entire image. This is due to the fact that NeRF needs to query hundreds of sample points along the ray to render a single pixel.

%Using our patch-based algorithm to select the patch $P(I_n)$
%from input image $I_n$ and its paired patch $P(\widetilde{I}_n)$ from the watermarked result $\widetilde{I}_n$, the content loss is computed between them. In detail, the content loss has two component. One is the pixel-wise MSE loss $L_{mse}$, which can be calculated as \Eref{eq:mse}. The other content is perceptual loss $L_{percep}$, which is defined as:
%\begin{equation}
%  L_{percep} = \| - \|_2^2,
%  \label{eq:percep}
%\end{equation}
%where XX denotes the feature representation obtained from the relu4\_1 layer of an ImageNet-pretrained VGG-19 network. 

%\subsection{Distortions.}
%\subsection{Message Extraction.}
% We adopt an ordinary NeRF \cite{10.1007/978-3-030-58452-8_24} as our geometry generator, which is trained to model the opacity field $\sigma(x)$ and original radiance color field $c(x, d)$. 
\begin{table*}[!htb]
%  \centering
%  \makeatletter\def\@captype{table}
 \caption{Bit accuracies with different lengths compared with baselines. The results are averaged on all all examples.}
 \label{tab:comparison1}
 \centering
 \begin{tabular}{l|ccccc}
    \toprule
         & 4 bits & 8 bits & 16 bits & 32 bits & 48 bits \\
    \midrule
    \textbf{Proposed \CopyNeRF} &\textbf{100\%} & \textbf{100\%}  & \textbf{91.16\%} & \textbf{78.08\%}  & \textbf{60.06\%}   \\
    HiDDeN~\cite{zhang2020udh}+NeRF\cite{nerf2020} & 50.31\% & 50.25\% & 50.19\% & 50.11\% & 50.04\%      \\
    MBRS~\cite{jia2021mbrs}+NeRF~\cite{nerf2020} & 53.25\% & 51.38\% & 50.53\% & 49.80\% & 50.14\%      \\
    NeRF\cite{nerf2020} with message & 72.50\%  & 63.19\% &52.22\%  & 50.00\% & 51.04\%\\
    \CopyNeRF in geometry &   76.75\% & 68.00\% &60.16\%  & 54.86\% & 53.36\%\\
    \bottomrule
\end{tabular}
\vspace{-5pt}
\label{tab:bavsdl}
\end{table*}


\subsection{Implementation details}
We implement our method using PyTorch. An eight-layer MLP with $256$ channels and the following two MLP branches are used to predict the original colors $\mathbf{c}$ and opacities $\sigma$, respectively. We train a ``coarse'' network along with a ``fine'' network network for importance sampling. we sample 32 points along each ray in the coarse model and 64 points in the fine model. Next, the patch size is set to $150\times 150$. The hyperparameters in~\Eref{eq:content} and~\Eref{eq:loss} are set as $\lambda_1=0.01$, $\gamma_1=1$, and $\gamma_2=10.00$. We use the Adam optimizer with defaults values $\beta_1 = 0.9$, $\beta_2 = 0.999$, $\epsilon = 10^{-8}$, and a learning rate $5\times10^{-4}$ that decays following the exponential scheduler during optimization. In our experiments, we set $N_m$ in \Eref{eq:messageembedding} as  $256$. We first optimize MLPs $\Theta_{\sigma}$ and $\Theta_{c}$ using loss function~\Eref{eq:recon} for $200$K and $100$K iterations for Blender dataset~\cite{nerf2020} and LLFF dataset~\cite{mildenhall2019local} separately, and then train the models $E_\xi$, $D_\phi$, and $H_\chi$ on a single NVIDIA V100 GPU. During training, we have considered messages with different bit lengths and forms. If a message has $4$ bits, we take into account all $2^{4}$ situations during training. The model creator can choose one message considered in our training as the desired message.

% Our transmission encoder can extract transmission features of multiple resolutions by leveraging advantages from ERRNet~\cite{wei2019single} and U-Net~\cite{ronneberger2015u}. Based on our experiments, even a simple U-Net~\cite{ronneberger2015u} can also help to resist the reflection interference during rendering. The transmission encoder is first integrated with a feature decoder to compress the information to a lower dimension, and the output is fed into the transmission MLP.

% For the pre-training of the transmission encoder, we follow the common way adopted by most image restoration methods, where the $\loss_1$ and adversarial loss between the estimation and its ground truth are employed to optimize the solution space. 

% The training procedure for our 3D scene watermarking
% approach contains the geometric training stage and the watermarking training stage. We describe these two stages and the corresponding objective functions in the following.


\section{Experiments}
\subsection{Experimental settings}
\noindent\textbf{Dataset.} To evaluate our methods, we train and test our model on Blender dataset~\cite{nerf2020} and LLFF dataset~\cite{mildenhall2019local}, which are common datasets used for NeRF. Blender dataset contains 8 detailed synthetic objects with 100 images taken from virtual cameras arranged on a hemisphere pointed inward. As in NeRF~\cite{nerf2020}, for each scene we input 100 views for training. LLFF dataset consists of 8 real-world scenes that contain mainly forward-facing images. Each scene contains 20 to 62 images. The data split for this dataset also follows NeRF~\cite{nerf2020}. For each scene, we select $20$ images from their testing dataset to evaluate the visual quality. For the evaluation of bit accuracy, we render $200$ views for each scenes to test whether the message can be effectively extracted under different viewpoints. We report average values across all testing viewpoints in our experiments.

\noindent\textbf{Baselines.} To the best of our knowledge, there is no method specifically for protecting the copyright of NeRF models. We, therefore, compare with four strategies to guarantee a fair comparison: 1) \textbf{HiDDeN~\cite{zhu2018hidden}+NeRF\cite{nerf2020}}: processing images with classical 2D watermarking method HiDDeN~\cite{zhu2018hidden} before training the NeRF model; 2) \textbf{MBRS~\cite{jia2021mbrs}+NeRF~\cite{nerf2020}}: processing images with state-of-the-art 2D watermarking method MBRS~\cite{jia2021mbrs} before training the NeRF model; 3) \textbf{NeRF with message}: concatenating the message $\mathbf{M}$ with location $\mathbf{x}$ and viewing direction $\mathbf{d}$ as the input of NeRF; 4) \textbf{\CopyNeRF in geometry}: changing our \CopyNeRF by fusing messages with the geometry to evaluate whether geometry is a good option for message embedding. %4) \textbf{\CopyNeRF with volume rendering}: We replace our distortion-resistant rendering with a volume rendering used by vanilla NeRF to see whether such a rendering scheme is a viable solution .

%4) \textbf{\CopyNeRF with volume rendering}: We replace our distortion-resistant rendering with a volume rendering used by vanilla NeRF to see whether such a rendering scheme is a viable solution.

% All the results in~\Sref{sec:experiments} are obtained using six views for training. Due to the page size limitation, the results of other methods trained under more views can be found in the supplementary material.2) UDH~\cite{zhang2020udh}+NeRF~\cite{nerf2020}: applying UDH networks to encode the multi-view images before training the NeRF model and retrieving the information by pretrianed reveal network; 

%For simplicity, we choose ERRNet in this place as the reflection removal method.
\noindent\textbf{Evaluation methodology.} We evaluate the performance of our proposed method against other methods by following the standard of digital watermarking about the invisibility, robustness, and capacity. For \textit{invisibility}, we evaluate the performance by using PSNR, SSIM, and LPIPS~\cite{zhang2018unreasonable} to compare the visual quality of the rendered results after message embedding. For \textit{robustness}, we will investigate whether the encoded messages can be extracted effectively by measuring the bit accuracy on different distortions. Besides normal situations, we consider the following distortions for message extraction: 1) Gaussian noise, 2) Rotation, 3) Scaling, and 4) Gaussian blur. For \textit{capacity}, following the setting in previous work for the watermarking of explicit 3D models~\cite{Yoo_2022_CVPR}, we investigate the invisibility and robustness under different message length as $N_b\in\{4, 8, 16, 32, 48\}$, which has been proven effective in protecting 3D models~\cite{Yoo_2022_CVPR}. Since we have included different viewpoints in our experiments for each scene, our evaluation can faithfully reflect whether the evaluated method can guarantee its robustness and consistency across viewpoints. 

%\renejie{We compare the performance by considering the invisibility, robustness, and capacity. We report quantitative performance using PSNR, SSIM, and LPIPS to compare the visual quality of the rendered results. This can fully evaluate the percentage of image degradation caused by embedded messages. Then, we use XXX to evaluate}  

%We report quantitative performance using PSNR, SSIM and LPIPS~\cite{zhang2018unreasonable}. For PSNR and SSIM, higher value indicates better performance. For LPIPS~\cite{zhang2018unreasonable}, lower value indicates better performance. By measuring image fidelity using high-level features, LPIPS~\cite{zhang2018unreasonable} can better match human judgements of image similarity. 



% The bit length is $16$ and the accuracies are the averages of 200 random patches from different viewing directions.



% \begin{table*}[!htb]
% %  \centering
% %  \makeatletter\def\@captype{table}
%  \caption{Bit accuracies and reconstruction qualities compared with our baselines. Higher PSNR and SSIM values denote better results ($\uparrow$), while lower LPIPS values denote better results ($\downarrow$). We show the results on $N_b = 16$  bits. The results are averaged on all all examples.}
%  \label{tab:comparison2}
%  \centering
%  \begin{tabular}{l|cccc}
%     \toprule
%          & Bit Acc$\uparrow$ & PSNR $\uparrow$& SSIM $\uparrow$& LPIPS $\downarrow$ \\
%     \midrule
%     \textbf{Proposed \CopyNeRF} &\textbf{91.16\%}& 26.29  & 0.910 & 0.038   \\
%     HiDDeN~\cite{zhu2018hidden}+NeRF\cite{nerf2020} & 50.19\% & 26.53 & 0.917 & 0.035     \\
%     MBRS~\cite{jia2021mbrs}+NeRF~\cite{nerf2020} & 50.53\% & \textbf{28.79} & \textbf{0.925} & \textbf{0.022}  \\
%     NeRF with message & 52.22\% & 22.33 & 0.773 & 0.108\\
%     \CopyNeRF in geometry & 60.16\% & 20.24 & 0.771  & 0.095 \\
%     \bottomrule
% \end{tabular}
% \label{tab:bavq}
% \end{table*}
\begin{figure*}
  \begin{minipage}{0.72\linewidth}
    \centering
    \tabcaption{Bit accuracies and reconstruction qualities compared with our baselines. $\uparrow$ ($\downarrow$) means higher (lower) is better. We show the results on $N_b = 16$  bits. The results are averaged on all all examples. The best performances are highlighted in \textbf{bold}.}
\label{tab:comparison2}
 \begin{tabular}{l|cccc}
    \toprule
         & Bit Acc$\uparrow$ & PSNR $\uparrow$& SSIM $\uparrow$& LPIPS $\downarrow$ \\
    \midrule
    \textbf{Proposed \CopyNeRF} &\textbf{91.16\%}& 26.29  & 0.910 & 0.038   \\
    HiDDeN~\cite{zhu2018hidden}+NeRF\cite{nerf2020} & 50.19\% & 26.53 & 0.917 & 0.035     \\
    MBRS~\cite{jia2021mbrs}+NeRF~\cite{nerf2020} & 50.53\% & \textbf{28.79} & \textbf{0.925} & \textbf{0.022}  \\
    NeRF with message & 52.22\% & 22.33 & 0.773 & 0.108\\
    \CopyNeRF in geometry & 60.16\% & 20.24 & 0.771  & 0.095 \\
    \bottomrule
\end{tabular}
\label{tab:bavq}
\end{minipage}
\begin{minipage}{0.28\linewidth}
    \centering
    \includegraphics[width = 0.76\linewidth]{figs/re_compare_1.pdf}
    \vspace{-5pt}
    \figcaption{Analysis for failure of MBRS~\cite{jia2021mbrs}+NeRF.}
    \label{fig:compare}
  \end{minipage}
\end{figure*}


\begin{table*}[!htb]
%  \centering
%  \makeatletter\def\@captype{table}
 \caption{Bit accuracies with different distortion types compared with each baseline and our \CopyNeRF without distortion-resistant rendering (DRR).  We show the results on $N_b = 16$ bits. The results are averaged on all all examples.}
 \label{tab:comparison3}
 \centering
 \begin{tabular}{l|ccccc}
    \toprule
         & No Distortion & Gaussian noise  & Rotation  & Scaling & Gaussian blur\\
         & & ($\nu$=0.1) & ($\pm\pi/6$) & ($\le25\%$) & ($deviation = 0.1$) \\
    \midrule
    \textbf{Proposed \CopyNeRF} &91.16\% & \textbf{90.44\%}  & \textbf{88.13\%} & \textbf{89.33\%} &  \textbf{90.06\%}   \\
    HiDDeN~\cite{zhu2018hidden}+NeRF\cite{nerf2020} & 50.19\% & 49.84\% & 50.12\% & 50.09\% &  50.16\%    \\
    MBRS~\cite{jia2021mbrs}+NeRF~\cite{nerf2020} &  50.53\% & 51.00\%  & 51.03\% & 50.12\% & 50.41\% \\
    NeRF with message & 52.22\%  & 50.53\%       & 50.22\%  & 50.19\% & 51.34\%\\
    \CopyNeRF in geometry &   60.16\% & 58.00\% &56.94\%  & 60.09\% & 59.38\%\\
    \CopyNeRF W/o DRR & \textbf{91.25\%}  & 89.12\%  & 75.81\%  & 87.44\% & 87.06\%\\
    \bottomrule
\end{tabular}
\vspace{-5pt}
 \label{tab:distortion}
 \end{table*}



\subsection{Experimental results}
\noindent\textbf{Qualitative results.}
\label{sec:qualitativeresults}
We first compare the reconstruction quality visually against all baselines and the results are shown in~\Fref{fig:visualquality}. Actually, all methods except NeRF with message and \CopyNeRF in geometry can achieve high reconstruction quality. For HiDDeN~\cite{zhu2018hidden} + NeRF~\cite{nerf2020} and MBRS~\cite{jia2021mbrs}+NeRF~\cite{nerf2020}, although they are efficient approaches in 2D watermarking, their bit accuracy values are all low for rendered images, which proves that the message are not effectively embedded after NeRF model training. From the results shown in \Fref{fig:compare}, the view synthesis of NeRF changes the information embedded by 2D watermarking methods, leading to their failures. For NeRF with message, as assumed in our previous discussions, directly employing secret messages as an input change the appearance of the output, which leads to their lower PSNR values. Besides, its lower bit accuracy also proves that this is not an effective embedding scheme. For \CopyNeRF in geometry, it achieves the worst visual quality among all methods. The rendered results look blurred, which confirms our assumption that the geometry is not a good option for message embedding. 



%\subsection{Quantitative results}
%\label{sec:5.2}
\noindent\textbf{Bit Accuracy vs. Message Length.}
We launch $5$ experiments for each message length and show
the relationship between bit accuracy and the length of message in~\Tref{tab:bavsdl}. We could clearly see that the bit accuracy drops when the number of bits increases. However, our \CopyNeRF achieves the best bit accuracy across all settings, which proves that the messages can be effectively embedded and robustly extracted. \CopyNeRF in geometry achieves the second best results among all setting, which shows that embedding message in geometry should also be a potential option for watermarking. However, the higher performance of our proposed \CopyNeRF shows that color representation is a better choice. 

% \vspace{-15pt}

\begin{figure*}
  \centering
  \includegraphics[width=\linewidth]{figs/Supp_diff_rendering.pdf}
  \caption{Comparisons for different rendering degradadtion in the inference phase. The message length is set to 16. We use average sampling points (ASP), importance sampling points (ISP), and random sampling points (RSP) in different rendering strategies. ``$32$ ASP + $32$ ISP'' is a strategy employed in the training process, and message extraction also shows the highest bit accuracy. When sampling strategies are changed to other ones during inference, the message extraction still shows similar performance, which verifies the effectiveness of our distortion-resistant rendering.}%Hidden~\cite{zhu2018hidden} + NeRF~\cite{nerf2020} can achieve higher PSNR, while its bit accuracy is very low. NeRF with messages can achieve higher bit accuracy when compared with Hidden~\cite{zhu2018hidden} + NeRF~\cite{nerf2020}, while its embedding scheme deteriorate its PSNR values.
  \label{fig:diff_render}
  \vspace{-12pt}
\end{figure*}

\noindent\textbf{Bit Accuracy vs. Reconstruction Quality.}
 We conduct more experiments to evaluate the relationship between bit accuracy and reconstruction quality. The results are shown in~\Tref{tab:bavq}\footnote{Results for other lengths of raw bits can be found in the supplementary materials.}. Our proposed \CopyNeRF achieves a good balance between bit accuracy and error metric values. Though the visual quality values are not the highest, the bit accuracy is the best among all settings. Though HiDDeN~\cite{zhu2018hidden} + NeRF~\cite{nerf2020} and MBRS~\cite{jia2021mbrs}+NeRF~\cite{nerf2020} can produce better visual quality values, its lower bit accuracy indicates that the secret messages are not effectively embedded and robustly extracted. NeRF with message also achieves degraded performance on bit accuracy, and its visual quality values are also low. It indicates that the embedded messages undermine the quality of reconstruction. Specifically, the lower visual quality values of \CopyNeRF in geometry indicates that hidding messages in color may lead to better reconstruction quality than hiding messages in geometry.


%In general, high bit accuracy results in higher PSRN, SSIM, and better LPIPS values. Specifically, HiDDen + NeRF achieves better PSNR, SSIM, and LPIPS values, since the invisible watermakred embedd into 2D images does not obviously change  


%Our method, Hidden + NeRF, and UDH + NeRF all achieves high error metric values, which are in line with the observation in~\Sref{sec:qualitativeresults}. Their consistently lower bit accuracy can be attributed to their poor performance in embedding and extracting messages. 

%Our method achieves better results across different message lengths, which shows that our method meets the requirements for a digital watermarking system: robustness, invisibility, and capacity. 


\noindent\textbf{Model robustness on 2D distortions.}
% We evaluate the robustness of our method to different distortions including noise, rotation, scaling and cropping in~\Tref{tab:distortion}. Specifically, we replace our distortion-resistant rendering with a volume rendering used by vanilla NeRF to see whether such a rendering scheme is a viable solution. We could see that our method is quite robust to different distortions. Specifically, \CopyNeRF w/o DRR achieves similar performance to the complete \CopyNeRF when no distortion is encountered. However, when it comes to different distortions, its lower bit accuracies demonstrates the effectiveness of our distortion-resistant rendering during training. 
We evaluate the robustness of our method by appyling several traditional 2D distortions. Specifically, as shown in~\Tref{tab:distortion}, we consider several types of 2D distortions including noise, rotation, scaling, and cropping. We could see that our method is quite robust to different 2D distortions. Specifically, \CopyNeRF w/o DRR achieves similar performance to the complete \CopyNeRF when no distortion is encountered. However, when it comes to different distortions, its lower bit accuracies demonstrates the effectiveness of our distortion-resistant rendering during training.

\noindent\textbf{Analysis for feature field.}
\label{sec:5.3}
In the section, we further evaluate the effectiveness of color feature field and message feature field. We first remove the module for building color feature field and directly combine the color representation with the message features. In this case, the model performs poorly in preserving the visual quality of the rendered results. We further remove the module for building message feature field and combine the message directly with the color feature field. The results in~\Tref{tab:ablationstudy} indicate that this may result in lower bit accuracy, which proves that messages are not embedded effectively. 

\begin{table}[tb]
%  \centering
%  \makeatletter\def\@captype{table}
 \caption{Comparisons for our full model, our model without Message Feature Field (MFF) and our model without Color Feature Field (CFF). The last row shows that our method achieves consistent performance even when different rendering scheme (DRS) is applied during \underline{testing}. }
 \label{tab:comparison4}
 \centering
 \begin{tabular}{l|cccc}
    \toprule
         &Bit Acc$\uparrow$& PSNR $\uparrow$& SSIM $\uparrow$& LPIPS $\downarrow$ \\
    \midrule
    \textbf{Ours} & \textbf{100\%} & \textbf{32.68} & \textbf{0.948} & \textbf{0.048}  \\
    W/o MFF& 82.69\% & 20.46 & 0.552 & 0.285     \\
    W/o CFF& 80.69\% & 21.06 & 0.612 & 0.187     \\
    DRS  & \textbf{100\%}  & 32.17 & 0.947 & 0.052   \\
    \bottomrule
\end{tabular}
\label{tab:ablationstudy}
      \vspace{-12pt}
 \end{table}
\noindent\textbf{Model robustness on rendering.}
% Though we use a distortion-resistant rendering strategy for training, the messages can also be effectively extracted using the volume rendering of NeRF. As shown in the last row of~\Tref{tab:ablationstudy}, the quantitative values with the volume rendering are still similar to original results in the first row of~\Tref{tab:ablationstudy}, which further confirms the robustness of our proposed method.
Though we apply a normal volume rendering strategy for inference, the messages can also be effectively extracted using the a distortion rendering utilized in training phase. As shown in the last row of~\Tref{tab:ablationstudy}, the quantitative values with the distortion rendering are still similar to original results in the first row of~\Tref{tab:ablationstudy}, which further confirms the robustness of our proposed method.

The results for different sampling schemes are presented in~\Fref{fig:diff_render}. Our distortion-resistant rendering employs $32$ average sampling points and $32$ importance sampling points during training. When different sampling strategies are applied in the inference phase, our method can also achieve high bit accuracy, which can validate the robustness of our method referring to different sampling strategies.


\noindent\textbf{Compasion with NeRF+HiDDeN/MBRS~\cite{zhu2018hidden, jia2021mbrs}.} We also conduct an experiment to compare the our method with approaches by directy applying 2D watermaring method on rendered images, namely NeRF+HiDDeN~\cite{zhu2018hidden} and NeRF+MBRS~\cite{jia2021mbrs}. Although these methods can reach a high bit accuracy as reported in their papers, 
as shown in \Fref{fig:hotdog}, these methods can easily leave detectable traces especially in areas with lower geometry values, as they lack the consideration for 3D information during watermarking. Besides, they only consider the media in 2D domain and cannot protect the NeRF model weights.
%We employ a random sampling strategy to ensure that our method is also robust even with different sampling strategy. One experiments are conducted in this section to show that even with different sampling strategy our method can also show competitive bit accuracy and visual quality values, while a module with deterministic sampling strategy shows degraded performances. 
%\begin{figure}
%  \centering
%  \includegraphics[width=\linewidth]{figs/fig_e3.pdf}
%  \caption{Visual quality comparison of the ablation study. We only show one frame for display purposes. From left to right: the ground truth, the result obtained by our full model, the model without message representation, and the model without color representation.}
%  \label{fig:ablationstudy}
%\end{figure}
\begin{figure}
  \centering
  \includegraphics[width=\linewidth]{figs/nerf_hidden_hotdog.pdf}
  \caption{Comparisons for watermarking after rendering. The patch in the lower left corner shows the augmentation result by simply multiplying a factor 30. We use image inversion for better visualization}
  \label{fig:hotdog}
  \vspace{-15pt}
\end{figure}
\section{Conclusions}
% In this paper, by ensuring any novel views synthesized from implicit scene representation are watermarked, our proposed framework can help to claim its ownership once shared such representation with others. In order to guarantee the invisibility of embedded messages, we keep the geometry and create a watermarked color representation to embed the message. Using color, position, viewing angle, and message to output a watermarked color, our method can guarantee the invisibility of copyright messages in the rendered results. Additionally, we introduce a distortion-resistant rendering scheme to ensure robust message extraction even with a CNN-based classical message extractor or different rendering strategy. This method achieves a promising balance between bit accuracy and high visual quality in experimental evaluations.

In this paper, we propose a framework to create a copyright-embedded 3D implicit representation by embedding messages into model weights. In order to guarantee the invisibility of embedded information, we keep the geometry unchanged and construct a watermarked color representation to produce the message embedded color. The embedded message can be extracted by a CNN-based extractor from rendered images from any viewpoints, while keeping high reconstruction quality. Additionally, we introduce a distortion-resistant rendering scheme to enhance the robustness of our model under different types of distortion, including classical 2D degradation and different rendering strategies. The proposed method achieves a promising balance between bit accuracy and high visual quality in experimental evaluations.

\noindent\textbf{Limitations.} Though our method has shown promising performance in claiming the ownership of Neural Radiance Fields,  
training a NeRF model is time-consuming. We will consider how to speed up the training process in our future work. Besides, though we have considered several designs to strengthen the system robustness, this standard may still be undermined when malicious users directly attack model weights, \textit{ie}, the model weights are corrupted. We conduct a simple experiment by directly adding Gaussian noise (std = 0.01) to the model parameters, and the accuracy slightly decreases to $93.97\%$ ($N_b = 8$). As this may also affect rendering quality, such model weights corruption may not be a priority for malicious users who intend to display the content. We will still actively consider how to handle such attacks in our future work. 
 
\noindent\textbf{Acknowledgement.}  Renjie Wan is supported by the Blue Sky Research Fund of HKBU under Grant No. BSRF/21-22/16 and Guangdong Basic and Applied Basic Research Foundation under Grant No. 2022A1515110692. Qing Guo is supported by the A*STAR Centre for Frontier AI Research and the National Research Foundation, Singapore, and DSO National Laboratories under the AI Singapore Programme (AISG Award No: AISG2-GC-2023-008).
% We thank the NVIDIA Academic Hardware Grant Program for its support of our work.

\linespread{1.0}
\normalem
{\small
\bibliographystyle{ieee_fullname}
\bibliography{egbib}
}

\end{document}
